\documentclass[12pt]{article}
\usepackage[utf8]{inputenc}
\usepackage[T1]{fontenc}
\usepackage{hyperref}
\title{Mahakam Kenyah Language Resources for K-12 Curriculum Development}
\author{Compiled by Claude Sonnet 5.5 (claude-sonnet-5-5) for Andr\'as Kornai}
\date{}
\begin{document}
\maketitle

\section{Introduction}
Mahakam Kenyah (ISO 639-3 xkm; Glottocode maha1308) is a Kenyah variety of East Kalimantan, Indonesia. The identification line in the task gave ``Bookkeeping'' as the family, which is clearly a data error and was not used. The search returned no page dedicated to this language, so this report is deliberately minimal.

\section{Primary Educational Resources}
None found.

\section{Linguistic and Reference Materials}
None specific to Mahakam Kenyah. Related Mainstream Kenyah background is in the Wikipedia overview: \url{https://en.wikipedia.org/wiki/Mainstream_Kenyah_language}. This is a different ISO language (xkl), so treat it as context only.

\section{Religious and Community Resources}
None found.

\section{Digital and Audio Resources}
Global Recordings Network pages for neighbouring Mainstream Kenyah (xkl) varieties exist, for example Lepo' Tau: \url{https://www.globalrecordings.net/en/language/4137}. They are not confirmed to be Mahakam Kenyah.

\section{Curriculum Development Context}
There is no evidence in the search results of an orthography, school programme or published materials for this language. Fieldwork and documentation would be needed first.
\end{document}